\section{Affine bundles}
\label{section-affine-vector}

\noindent
For an affine bundle the pullback map is surjective on Chow groups.

\begin{lemma}
\label{lemma-pullback-affine-fibres-surjective}
Let $(S, \delta)$ be as in Situation \ref{situation-setup}.
Let $X$, $Y$ be locally of finite type over $S$.
Let $f : X \to Y$ be a flat morphism of relative dimension $r$.
Assume that for every $y \in Y$, there exists an open neighbourhood
$U \subset Y$ such that $f|_{f^{-1}(U)} : f^{-1}(U) \to U$
is identified with the morphism $U \times \mathbf{A}^r \to U$.
Then $f^* : \CH_k(Y) \to \CH_{k + r}(X)$ is surjective for all
$k \in \mathbf{Z}$.
\end{lemma}

\begin{proof}
Let $\alpha \in \CH_{k + r}(X)$.
Write $\alpha = \sum m_j[W_j]$ with $m_j \not = 0$ and
$W_j$ pairwise distinct integral closed subschemes of
$\delta$-dimension $k + r$. Then the family $\{W_j\}$
is locally finite in $X$. For any quasi-compact open
$V \subset Y$ we see that $f^{-1}(V) \cap W_j$
is nonempty only for finitely many $j$. Hence the
collection $Z_j = \overline{f(W_j)}$ of closures
of images is a locally finite collection of integral
closed subschemes of $Y$.

\medskip\noindent
Consider the fibre product diagrams
$$
\xymatrix{
f^{-1}(Z_j) \ar[r] \ar[d]_{f_j} & X \ar[d]^f \\
Z_j \ar[r] & Y
}
$$
Suppose that $[W_j] \in Z_{k + r}(f^{-1}(Z_j))$
is rationally equivalent to $f_j^*\beta_j$ for some
$k$-cycle $\beta_j \in \CH_k(Z_j)$. Then
$\beta = \sum m_j \beta_j$ will be a $k$-cycle on $Y$
and $f^*\beta = \sum m_j f_j^*\beta_j$ will be rationally
equivalent to $\alpha$ (see
Remark \ref{remark-infinite-sums-rational-equivalences}).
This reduces us to the case $Y$ integral, and
$\alpha = [W]$ for some integral closed subscheme
of $X$ dominating $Y$. In particular we may
assume that $d = \dim_\delta(Y) < \infty$.

\medskip\noindent
Hence we can use induction on $d = \dim_\delta(Y)$.
If $d < k$, then $\CH_{k + r}(X) = 0$ and the lemma holds.
By assumption there exists a dense open $V \subset Y$ such
that $f^{-1}(V) \cong V \times \mathbf{A}^r$ as schemes over $V$.
Suppose that we can show that $\alpha|_{f^{-1}(V)} = f^*\beta$
for some $\beta \in Z_k(V)$. By Lemma \ref{lemma-exact-sequence-open}
we see that
$\beta = \beta'|_V$ for some $\beta' \in Z_k(Y)$.
By the exact sequence
$\CH_k(f^{-1}(Y \setminus V)) \to \CH_k(X) \to \CH_k(f^{-1}(V))$
of Lemma \ref{lemma-restrict-to-open}
we see that $\alpha - f^*\beta'$ comes from
a cycle $\alpha' \in \CH_{k + r}(f^{-1}(Y \setminus V))$.
Since $\dim_\delta(Y \setminus V) < d$ we win by
induction on $d$.

\medskip\noindent
Thus we may assume that $X = Y \times \mathbf{A}^r$.
In this case we can factor $f$ as
$$
X = Y \times \mathbf{A}^r \to
Y \times \mathbf{A}^{r - 1} \to \ldots \to
Y \times \mathbf{A}^1 \to Y.
$$
Hence it suffices to do the case $r = 1$. By the argument in the
second paragraph of the proof we are reduced to the case
$\alpha = [W]$, $Y$ integral, and $W \to Y$ dominant.
Again we can do induction on $d = \dim_\delta(Y)$.
If $W = Y \times \mathbf{A}^1$, then $[W] = f^*[Y]$.
Lastly, $W \subset Y \times \mathbf{A}^1$ is a proper inclusion,
then $W \to Y$ induces a finite field extension $R(W)/R(Y)$.
Let $P(T) \in R(Y)[T]$ be the monic irreducible polynomial such
that the generic fibre of $W \to Y$ is cut out by $P$ in
$\mathbf{A}^1_{R(Y)}$. Let $V \subset Y$ be a nonempty open such
that $P \in \Gamma(V, \mathcal{O}_Y)[T]$, and such that
$W \cap f^{-1}(V)$ is still cut out by $P$. Then we see that
$\alpha|_{f^{-1}(V)} \sim_{rat} 0$ and hence $\alpha \sim_{rat} \alpha'$
for some cycle $\alpha'$ on $(Y \setminus V) \times \mathbf{A}^1$.
By induction on the dimension we win.
\end{proof}

\begin{lemma}
\label{lemma-linebundle}
Let $(S, \delta)$ be as in Situation \ref{situation-setup}.
Let $X$ be locally of finite type over $S$.
Let $\mathcal{L}$ be an invertible $\mathcal{O}_X$-module.
Let
$$
p :
L = \underline{\Spec}(\text{Sym}^*(\mathcal{L}))
\longrightarrow
X
$$
be the associated vector bundle over $X$.
Then $p^* : \CH_k(X) \to \CH_{k + 1}(L)$ is an isomorphism for all $k$.
\end{lemma}

\begin{proof}
For surjectivity see Lemma \ref{lemma-pullback-affine-fibres-surjective}.
Let $o : X \to L$ be the zero section of $L \to X$, i.e., the morphism
corresponding to the surjection $\text{Sym}^*(\mathcal{L}) \to \mathcal{O}_X$
which maps $\mathcal{L}^{\otimes n}$ to zero for all $n > 0$.
Then $p \circ o = \text{id}_X$ and $o(X)$ is an effective
Cartier divisor on $L$. Hence by Lemma \ref{lemma-relative-effective-cartier}
we see that $o^* \circ p^* = \text{id}$ and we conclude that $p^*$ is
injective too.
\end{proof}

\begin{remark}
\label{remark-when-isomorphism}
We will see later (Lemma \ref{lemma-vectorbundle}) that if $X$ is a
vector bundle of rank $r$ over $Y$ then the pullback map
$\CH_k(Y) \to \CH_{k + r}(X)$
is an isomorphism. This is true whenever $X \to Y$ satisfies
the assumptions of Lemma \ref{lemma-pullback-affine-fibres-surjective}, see
\cite[Lemma 2.2]{Totaro-group}. We will sketch a proof in
Remark \ref{remark-higher-chow-isomorphism} using higher chow groups.
\end{remark}

\begin{lemma}
\label{lemma-linebundle-formulae}
In the situation of Lemma \ref{lemma-linebundle} denote $o : X \to L$
the zero section (see proof of the lemma). Then we have
\begin{enumerate}
\item $o(X)$ is the zero scheme of a regular global section of
$p^*\mathcal{L}^{\otimes -1}$,
\item $o_* : \CH_k(X) \to \CH_k(L)$ as $o$ is a closed immersion,
\item $o^* : \CH_{k + 1}(L) \to \CH_k(X)$ as $o(X)$
is an effective Cartier divisor,
\item $o^* p^* : \CH_k(X) \to \CH_k(X)$ is the identity map,
\item $o_*\alpha = - p^*(c_1(\mathcal{L}) \cap \alpha)$ for any
$\alpha \in \CH_k(X)$, and
\item $o^* o_* : \CH_k(X) \to \CH_{k - 1}(X)$ is equal to the map
$\alpha \mapsto - c_1(\mathcal{L}) \cap \alpha$.
\end{enumerate}
\end{lemma}

\begin{proof}
Since $p_*\mathcal{O}_L = \text{Sym}^*(\mathcal{L})$ we have
$p_*(p^*\mathcal{L}^{\otimes -1}) =
\text{Sym}^*(\mathcal{L}) \otimes_{\mathcal{O}_X} \mathcal{L}^{\otimes -1}$
by the projection formula
(Cohomology, Lemma \ref{cohomology-lemma-projection-formula})
and the section mentioned in (1) is
the canonical trivialization $\mathcal{O}_X \to
\mathcal{L} \otimes_{\mathcal{O}_X} \mathcal{L}^{\otimes -1}$.
We omit the proof that the vanishing locus of
this section is precisely $o(X)$. This proves (1).

\medskip\noindent
Parts (2), (3), and (4) we've seen in the course of the proof of
Lemma \ref{lemma-linebundle}. Of course (4) is the first
formula in Lemma \ref{lemma-relative-effective-cartier}.

\medskip\noindent
Part (5) follows from the second formula in
Lemma \ref{lemma-relative-effective-cartier},
additivity of capping with $c_1$ (Lemma \ref{lemma-c1-cap-additive}),
and the fact that capping with $c_1$ commutes with flat pullback
(Lemma \ref{lemma-flat-pullback-cap-c1}).

\medskip\noindent
Part (6) follows from Lemma \ref{lemma-gysin-back}
and the fact that $o^*p^*\mathcal{L} = \mathcal{L}$.
\end{proof}

\begin{lemma}
\label{lemma-decompose-section}
Let $Y$ be a scheme. Let $\mathcal{L}_i$, $i = 1, 2$ be invertible
$\mathcal{O}_Y$-modules. Let $s$ be a global section of
$\mathcal{L}_1 \otimes_{\mathcal{O}_Y} \mathcal{L}_2$.
Denote $i : D \to Y$ the zero scheme of $s$.
Then there exists a commutative diagram
$$
\xymatrix{
D_1 \ar[r]_{i_1} \ar[d]_{p_1} &
L \ar[d]^p &
D_2 \ar[l]^{i_2} \ar[d]^{p_2} \\
D \ar[r]^i &
Y &
D \ar[l]_i
}
$$
and sections $s_i$ of $p^*\mathcal{L}_i$ such that
the following hold:
\begin{enumerate}
\item $p^*s = s_1 \otimes s_2$,
\item $p$ is of finite type and flat of relative dimension $1$,
\item $D_i$ is the zero scheme of $s_i$,
\item $D_i \cong
\underline{\Spec}(\text{Sym}^*(\mathcal{L}_{3 - i}^{\otimes -1})|_D))$
over $D$ for $i = 1, 2$,
\item $p^{-1}D = D_1 \cup D_2$ (scheme theoretic union),
\item $D_1 \cap D_2$ (scheme theoretic intersection) maps
isomorphically to $D$, and
\item $D_1 \cap D_2 \to D_i$
is the zero section of the line bundle $D_i \to D$ for $i = 1, 2$.
\end{enumerate}
Moreover, the formation of this diagram and the sections $s_i$
commutes with arbitrary base change.
\end{lemma}

\begin{proof}
Let $p : X \to Y$ be the relative spectrum of the quasi-coherent
sheaf of $\mathcal{O}_Y$-algebras
$$
\mathcal{A} =
\left(\bigoplus\nolimits_{a_1, a_2 \geq 0}
\mathcal{L}_1^{\otimes -a_1} \otimes_{\mathcal{O}_Y}
\mathcal{L}_2^{\otimes -a_2}\right)/\mathcal{J}
$$
where $\mathcal{J}$ is the ideal generated by local sections of
the form $st - t$ for $t$ a local section of any summand
$\mathcal{L}_1^{\otimes -a_1} \otimes \mathcal{L}_2^{\otimes -a_2}$
with $a_1, a_2 > 0$. The sections $s_i$ viewed as maps
$p^*\mathcal{L}_i^{\otimes -1} \to \mathcal{O}_X$ are defined as the adjoints
of the maps $\mathcal{L}_i^{\otimes -1} \to \mathcal{A} = p_*\mathcal{O}_X$.
For any $y \in Y$ we can choose an affine
open $V \subset Y$, say $V = \Spec(A)$, containing $y$ and
trivializations $z_i : \mathcal{O}_V \to \mathcal{L}_i^{\otimes -1}|_V$.
Observe that $f = s(z_1z_2) \in A$ cuts out the closed subscheme $D$.
Then clearly
$$
p^{-1}(V) = \Spec(B[z_1, z_2]/(z_1 z_2 - f))
$$
Since $D_i$ is cut out by $z_i$ everything is clear.
\end{proof}

\begin{lemma}
\label{lemma-decompose-section-formulae}
In the situation of Lemma \ref{lemma-decompose-section}
assume $Y$ is locally of finite type over $(S, \delta)$ as in
Situation \ref{situation-setup}. Then we have
$i_1^*p^*\alpha = p_1^*i^*\alpha$
in $\CH_k(D_1)$ for all $\alpha \in \CH_k(Y)$.
\end{lemma}

\begin{proof}
Let $W \subset Y$ be an integral closed subscheme of $\delta$-dimension $k$.
We distinguish two cases.

\medskip\noindent
Assume $W \subset D$. Then
$i^*[W] = c_1(\mathcal{L}_1) \cap [W] + c_1(\mathcal{L}_2) \cap [W]$
in $\CH_{k - 1}(D)$ by our definition of Gysin homomorphisms and the
additivity of Lemma \ref{lemma-c1-cap-additive}.
Hence $p_1^*i^*[W] =
p_1^*(c_1(\mathcal{L}_1) \cap [W]) + p_1^*(c_1(\mathcal{L}_2) \cap [W])$.
On the other hand, we have
$p^*[W] = [p^{-1}(W)]_{k + 1}$ by construction of flat pullback.
And $p^{-1}(W) = W_1 \cup W_2$ (scheme theoretically)
where $W_i = p_i^{-1}(W)$ is a line bundle over $W$
by the lemma (since formation of the diagram commutes with base change).
Then $[p^{-1}(W)]_{k + 1} = [W_1] + [W_2]$ as $W_i$ are integral closed
subschemes of $L$ of $\delta$-dimension $k + 1$. Hence
\begin{align*}
i_1^*p^*[W]
& =
i_1^*[p^{-1}(W)]_{k + 1} \\
& =
i_1^*([W_1] + [W_2]) \\
& =
c_1(p_1^*\mathcal{L}_1) \cap [W_1] + [W_1 \cap W_2]_k \\
& =
c_1(p_1^*\mathcal{L}_1) \cap p_1^*[W] + [W_1 \cap W_2]_k \\
& =
p_1^*(c_1(\mathcal{L}_1) \cap [W]) + [W_1 \cap W_2]_k
\end{align*}
by construction of Gysin homomorphisms, the definition of flat pullback
(for the second equality), and compatibility of $c_1 \cap -$
with flat pullback (Lemma \ref{lemma-flat-pullback-cap-c1}).
Since $W_1 \cap W_2$ is the zero section of the line bundle
$W_1 \to W$ we see from Lemma \ref{lemma-linebundle-formulae}
that $[W_1 \cap W_2]_k = p_1^*(c_1(\mathcal{L}_2) \cap [W])$.
Note that here we use the fact that $D_1$ is the line bundle
which is the relative spectrum of the inverse of $\mathcal{L}_2$.
Thus we get the same thing as before.

\medskip\noindent
Assume $W \not \subset D$. In this case, both $i_1^*p^*[W]$
and $p_1^*i^*[W]$ are represented by the $k - 1$ cycle associated
to the scheme theoretic inverse image of $W$ in $D_1$.
\end{proof}

\begin{lemma}
\label{lemma-normal-cone-effective-Cartier}
In Situation \ref{situation-setup} let $X$ be a scheme locally
of finite type over $S$. Let $(\mathcal{L}, s, i : D \to X)$
be a triple as in Definition \ref{definition-gysin-homomorphism}.
There exists a commutative diagram
$$
\xymatrix{
D' \ar[r]_{i'} \ar[d]_p & X' \ar[d]^g \\
D \ar[r]^i & X
}
$$
such that
\begin{enumerate}
\item $p$ and $g$ are of finite type and flat of relative dimension $1$,
\item $p^* : \CH_k(D) \to \CH_{k + 1}(D')$ is injective for all $k$,
\item $D' \subset X'$ is the zero scheme of a global section
$s' \in \Gamma(X', \mathcal{O}_{X'})$,
\item $p^*i^* = (i')^*g^*$ as maps $\CH_k(X) \to \CH_k(D')$.
\end{enumerate}
Moreover, these properties remain true after arbitrary base change
by morphisms $Y \to X$ which are locally of finite type.
\end{lemma}

\begin{proof}
Observe that $(i')^*$ is defined because we have the triple
$(\mathcal{O}_{X'}, s', i' : D' \to X')$ as in
Definition \ref{definition-gysin-homomorphism}. Thus the statement makes sense.

\medskip\noindent
Set $\mathcal{L}_1 = \mathcal{O}_X$, $\mathcal{L}_2 = \mathcal{L}$
and apply Lemma \ref{lemma-decompose-section} with the section $s$ of
$\mathcal{L} = \mathcal{L}_1 \otimes_{\mathcal{O}_X} \mathcal{L}_2$.
Take $D' = D_1$. The results now follow from the lemma, from
Lemma \ref{lemma-decompose-section-formulae}
and injectivity by
Lemma \ref{lemma-linebundle}.
\end{proof}

\begin{remark}
\label{remark-higher-chow-isomorphism}
Let $(S, \delta)$ be as in Situation \ref{situation-setup}.
Let $Y$ be locally of finite type over $S$. Let $r \geq 0$. Let
$f : X \to Y$ be a morphism of schemes. Assume every $y \in Y$
is contained in an open $V \subset Y$ such that
$f^{-1}(V) \cong V \times \mathbf{A}^r$ as schemes over $V$.
In this remark we sketch a proof of the fact that
$f^* : \CH_k(Y) \to \CH_{k + r}(X)$ is an isomorphism.
First, by Lemma \ref{lemma-pullback-affine-fibres-surjective}
the map is surjective. Let $\alpha \in \CH_k(Y)$ with $f^*\alpha = 0$.
We will prove that $\alpha = 0$.

\medskip\noindent
Step 1. We may assume that $\dim_\delta(Y) < \infty$. (This is immediate
in all cases in practice so we suggest the reader skip this step.)
Namely, any rational equivalence witnessing that $f^*\alpha = 0$ on $X$,
will use a locally finite collection of integral closed subschemes 
of dimension $k + r + 1$. Taking the union of the closures of the images
of these in $Y$ we get a closed subscheme $Y' \subset Y$
of $\dim_\delta(Y') \leq k + r + 1$ such that
$\alpha$ is the image of some $\alpha' \in \CH_k(Y')$
and such that $(f')^*\alpha = 0$ where $f'$ is the base change of $f$ to $Y'$.

\medskip\noindent
Step 2. Assume $d = \dim_\delta(Y) < \infty$. Then we can use induction on
$d$. If $d < k$, then $\alpha = 0$ and we are done; this is the base case
of the induction.
In general, our assumption on $f$ shows we can choose a dense open
$V \subset Y$ such that $U = f^{-1}(V) = \mathbf{A}^r_V$.
Denote $Y' \subset Y$ the complement of $V$ as a reduced closed
subscheme and set $X' = f^{-1}(Y')$. Consider
$$
\xymatrix{
\CH^M_{k + r}(U, 1) \ar[r] &
\CH_{k + r}(X') \ar[r] &
\CH_{k + r}(X) \ar[r] &
\CH_{k + r}(U) \ar[r] &
0 \\
\CH^M_k(V, 1) \ar[r] \ar[u] &
\CH_k(Y') \ar[r] \ar[u] &
\CH_k(Y) \ar[r] \ar[u] &
\CH_k(V) \ar[r] \ar[u] &
0
}
$$
Here we use the first higher Chow groups of $V$ and $U$ and the
six term exact sequences constructed in Remark \ref{remark-higher-chow},
as well as flat pullback for these higher chow groups and
compatibility of flat pullback with these six term exact sequences.
Since $U = \mathbf{A}^r_V$ the vertical map on the right is an isomorphism.
The map $\CH_k(Y') \to \CH_{k + r}(X')$ is bijective by induction on $d$.
Hence to finish the argument is suffices to show that
$$
\CH^M_k(V, 1) \longrightarrow \CH^M_{k + r}(U, 1)
$$
is surjective. Arguing as in the proof of
Lemma \ref{lemma-pullback-affine-fibres-surjective}
this reduces to Step 3 below.

\medskip\noindent
Step 3. Let $F$ be a field. Then $\CH^M_0(\mathbf{A}^1_F, 1) = 0$.
(In the proof of the lemma cited above we proved analogously that
$\CH_0(\mathbf{A}^1_F) = 0$.) We have
$$
\CH^M_0(\mathbf{A}^1_F, 1) = \Coker\left(
\partial : K^M_2(F(T)) \longrightarrow
\bigoplus\nolimits_{\mathfrak p \subset F[T]\text{ maximal}}
\kappa(\mathfrak p)^*\right)
$$
The classical argument for the vanishing of the cokernel is to
show by induction on the degree of $\kappa(\mathfrak p)/F$
that the summand corresponding to $\mathfrak p$ is in the image.
If $\mathfrak p$ is generated by the
irreducible monic polynomial $P(T) \in F[T]$ and if
$u \in \kappa(x)^*$ is the residue class of some
$Q(T) \in F[T]$ with $\deg(Q) < \deg(P)$ then one shows that
$\partial(Q, P)$ produces the element $u$ at $\mathfrak p$
and perhaps some other units at primes dividing $Q$ which
have lower degree. This finishes the sketch of the proof.
\end{remark}











